% docs/texte/05_theory/uv_coverage.tex -- theory chapter "Earth Rotation Synthesis and uv-Coverage" (Interferosim)
% Version: 0.1.0 (2026-10-12)
% Notation as defined in texte/05_theory/notation.tex.
% ================================================================

\index{uv-coverage}\index{Earth rotation synthesis}
\section{Baseline coordinates}

Baselines are expressed in an equatorial system: $X$ towards
$(H=0,\delta=0)$, $Y$ towards $(H=-6^{\mathrm h},\delta=0)$ (east) and
$Z$ towards the celestial pole. A baseline given in local coordinates
(east $E$, north $N$, up $U$) at latitude $\varphi$ becomes
\begin{equation}
  L_X = -N\sin\varphi + U\cos\varphi,\qquad
  L_Y = E,\qquad
  L_Z = N\cos\varphi + U\sin\varphi .
\end{equation}

\section{Spatial frequencies}

For a source at hour angle $H$ and declination $\delta$:
\begin{equation}
  \begin{pmatrix} u\\ v\\ w \end{pmatrix}
  = \frac{1}{\lambda}
  \begin{pmatrix}
    \sin H & \cos H & 0\\
    -\sin\delta\cos H & \sin\delta\sin H & \cos\delta\\
    \cos\delta\cos H & -\cos\delta\sin H & \sin\delta
  \end{pmatrix}
  \begin{pmatrix} L_X\\ L_Y\\ L_Z \end{pmatrix} .
  \label{eq:uvw}
\end{equation}
For an east--west baseline ($L_X=L_Z=0$, $L_Y=b$) this gives
\begin{equation}
  u = \frac{b}{\lambda}\cos H,\qquad
  v = \frac{b}{\lambda}\sin\delta\,\sin H .
\end{equation}

\section{Earth rotation synthesis}

As $H$ changes, each baseline traces an ellipse in the $(u,v)$ plane,
\begin{equation}
  u^2 + \left(\frac{v - (L_Z/\lambda)\cos\delta}{\sin\delta}\right)^2
  = \frac{L_X^2+L_Y^2}{\lambda^2},
\end{equation}
together with its Hermitian mirror image. The axis ratio is
$\sin\delta$; a source near the equator gives a strongly flattened track.
The usable hour-angle range is limited by the elevation of the source,
\begin{equation}
  \sin\varepsilon = \sin\varphi\sin\delta + \cos\varphi\cos\delta\cos H ,
\end{equation}
and a minimum elevation (Interferosim:
\texttt{Observation.min\_elev\_deg}).

\section{Coverage and image quality}

The set of measured $(u,v)$ points defines the sampling function
$\mathcal S(u,v)$. The longest projected baseline sets the resolution,
$\theta \approx \lambda/b_{\max}$; the shortest one sets the largest
angular scale that can be recovered, $\approx\lambda/b_{\min}$.
Emission on larger scales is filtered out. With a single baseline the
coverage is one ellipse; for the 11.5\,km example baseline less than
0.1\,\% of the cells of a typical image grid are occupied
(notebook~A).
